\subsection{Notation and basic conventions}
\label{sec:notation}

Vectors are columns unless they appear as rows of a data matrix; $v^\top$
denotes transpose. For a vector $v$, $\|v\|_2$ is its Euclidean norm and
$\|v\|_1=\sum_\ell|v_\ell|$. For a matrix $M$,
$\|M\|_{\rm F}=(\sum_{k,\ell}M_{k\ell}^2)^{1/2}$ is its Frobenius norm
and $\|M\|_{\rm op}=\sup_{\|v\|_2=1}\|Mv\|_2$ is its operator norm.
We use $\langle u,v\rangle=u^\top v$ and
$\langle M,N\rangle_{\rm F}=\operatorname{tr}(M^\top N)$ for vector and
matrix inner products; $\operatorname{tr}$ sums diagonal entries.
The symbols $\mathbf1_N$, $I_N$, and $e_\ell$ denote the all-ones vector,
identity matrix, and $\ell$th coordinate vector, respectively, with the
ambient dimension specified by context.

For a real number $z$, $z_+=\max\{z,0\}$ is its positive part.
Logarithms are natural; $\lfloor z\rfloor$ is the largest integer not exceeding $z$. We write $\mathbb E$ and $\Pr$ for expectation
and probability under the specified sampling law. For a square-integrable
random scalar or vector $V$, $\|V\|_{L^2}=(\mathbb E\|V\|_2^2)^{1/2}$;
for scalars, the norm inside the expectation is absolute value.
An overbar denotes a sample average, and a hat denotes a fitted or
estimated quantity. The projection $\Pi_{\mathcal Y}(v)$ is the closest
point to $v$ in a nonempty closed convex set $\mathcal Y$ in Euclidean
distance. For a box, this operation clips each coordinate to its allowed
interval. All random variables and prediction maps are measurable.

Table~\ref{tab:notation} collects the main quantities. The forecasting
construction and the statistical statements below give their precise domains
and sampling assumptions.

\begin{table}[t]
\centering\small
\renewcommand{\arraystretch}{1.12}
\begin{tabular}{@{}p{0.26\linewidth}p{0.69\linewidth}@{}}
\toprule
Symbol & Meaning \\
\midrule
$y_t$, $d_y$ & Observation at time $t$ and its dimension. \\
$x_t$, $m$ & Causal history and number of observed time steps in it. \\
$\mathcal H$, $H$, $q$ & Set of forecast horizons, its size, and target dimension; for a joint forecast, $q=d_yH$. \\
$\Phi_b$, $C_b$ & Feature block $b$ and its number of coordinates. \\
$\Gamma_s$, $p_s$ & Stack of the first $s$ feature blocks and total dimension $p_s=\sum_{b=1}^s C_b$. \\
$\kappa$, $\lambda$ & Decoder sensitivity budget and ridge regularization parameter. \\
$R(f)$, $R^*$ & Population mean squared error of a fitted forecast and infimum over a stated decoder class. \\
$W_{\kappa,A}$ & Population pairwise obstruction, using budget $\kappa$ and conditioning variable $A$; distinct from coefficient matrices $W_b$. \\
$n$, $K$ & Number of calibration episodes and number $K=\lfloor n/2\rfloor$ of disjoint episode pairs. \\
$Z_{i,jh}$, $Q_{r,jh}$ & Loss of candidate $j$ at horizon $h$ on episode $i$, and obstruction score on pair $r$. \\
$L_{jh}$, $U_{jh}$ & Lower obstruction endpoint and upper fitted-risk endpoint; the unindexed $U=\Gamma(X)$ instead denotes a random representation. \\
$D_h$, $B_h$ & Diameter of the known horizon-$h$ output set and squared-loss bound $B_h=D_h^2$. \\
$c(X)$, $\pi_a$, $N_a$, $K_a$ & Frozen stratum label, population stratum mass, sample count, and within-stratum pair count. \\
$W_{\rm str}$, $L_{\rm str}$, $\mathcal P$ & Stratified population witness, its lower endpoint, and the confidence set for stratum masses. \\
$c_j$, $c_-$, $c_+$, $c^*$ & Declared class cost, lower and upper certified cost endpoints, and minimum feasible cost. \\
$\epsilon_h$, $\delta$, $\alpha$ & Requested RMSE tolerance, total failure probability, and failure probability allocated to one confidence statement. \\
\bottomrule
\end{tabular}
\caption{Recurring notation. Mean squared error (MSE) is an expected squared Euclidean error; root-mean-square error (RMSE) is its square root. Experiment-specific coordinate averaging is stated separately.}
\label{tab:notation}
\end{table}
